\section{Open Research Artifact and Continuous Verification}
\label{sec:software-artifact}

The analysis is maintained as a public research-software artifact at
\begin{center}
\href{https://github.com/redxe/quantum-reuse-security}
{\texttt{github.com/redxe/quantum-reuse-security}}.
\end{center}
The repository's version-0.6.0 baseline separates the command-line interface,
circuit helpers, measurement utilities, metrics, state preparation, validation,
and analysis orchestration under \texttt{src/quantum\_reuse/}. The exact NumPy
backend remains the reference implementation because its statevector, branch,
and partial-trace operations are visible without depending on an SDK-specific
execution model.

\subsection{Reproduction commands}

A clean installation can run the complete analysis and fixed-input summary with
\begin{lstlisting}[style=python]
python -m quantum_reuse analyze --output run_output
python -m quantum_reuse fixed-input-summary
pytest -q
\end{lstlisting}
The analysis emits branch-conditioned CSV data, averaged fifth-wire states,
distinguishability metrics, information summaries, numerical-validation data,
a run report, and plots. The committed baseline CSV files are generated from
the same command used in continuous integration.

\subsection{Continuous-integration contract}

The GitHub Actions workflow performs four layers of checking:
\begin{enumerate}
  \item install and test on Python 3.9, 3.11, and 3.13;
  \item run Black formatting checks on the modular package surfaces, tests, and
  examples;
  \item run Flake8 and the Pytest regression suite; and
  \item rerun the deterministic analysis on Python 3.11 and compare the three
  principal CSV outputs with committed baselines using zero relative tolerance
  and absolute tolerance $10^{-12}$.
\end{enumerate}
The workflow uses read-only repository permissions and disables persisted
checkout credentials. Passing runs therefore establish that the documented
software environment reproduces the committed numerical tables across the
matrix; they do not by themselves establish hardware behavior or security
against an adversarial implementation.

\subsection{Public roadmap}

The remaining work is recorded as public, reviewable issues:
\begin{itemize}
  \item \href{https://github.com/redxe/quantum-reuse-security/issues/1}
  {Issue 1: reconstruct the exact educational detector acceptance register};
  \item \href{https://github.com/redxe/quantum-reuse-security/issues/2}
  {Issue 2: implement branch-level Qiskit parity with the NumPy model}; and
  \item \href{https://github.com/redxe/quantum-reuse-security/issues/3}
  {Issue 3: decompose the legacy monolithic analysis core while preserving the
  public API and deterministic outputs}.
\end{itemize}
This issue structure makes the unresolved scientific and engineering boundaries
visible to reviewers rather than hiding them behind a polished final result.
